\section{Background and Related Work}
\label{sec:background-related}

\subsection{Device Fingerprints as Authentication Context}

Browser fingerprinting combines software, hardware, and runtime attributes exposed through HTTP headers and Web APIs into a probabilistic representation of a client environment~\cite{laperdrix2020browser,eckersley2010unique}. The same mechanism supports several distinct applications: stateless tracking, fraud analysis, bot detection, account protection, and adaptive authentication. These applications must not be conflated. A field that is useful for population-level linkability may be unstable across software updates; an attribute that is inexpensive to collect may be inexpensive to spoof; and a fingerprint that distinguishes desktop configurations may collapse many mobile devices into similar profiles.

Alaca and van Oorschot analyze device fingerprinting specifically as an auxiliary Web-authentication factor. Rather than ranking methods by uniqueness alone, they classify 29 approaches and compare properties including stability, repeatability, passiveness, resource cost, spoofing difficulty, and distinguishability~\cite{alaca2016device}. This framing is directly relevant to HybridGuard: the goal is not to maximize the entropy of a single identifier but to collect low-friction context whose limitations are explicit. Spooren et al. provide an important counterweight, finding that mobile device fingerprints can exhibit substantial similarity even across models and brands, which weakens their reliability as a standalone trusted-device signal for RBA~\cite{spooren2015harmful}. Empirical work on deployed RBA also shows that real services combine multiple context factors and apply different thresholds rather than treating a fingerprint as a cryptographic identity~\cite{wiefling2019really}.

HybridGuard follows the conservative implication of this literature. It does not claim that \appfp identifies a physical device uniquely or that a relation alert proves account fraud. Instead, it treats multiple surfaces from the same controlled session as evidence about environment coherence. The downstream consumer may be an RBA policy, but the evaluated task in this paper is device-environment manipulation and cross-layer relation violation.

\subsection{Browser Fingerprint Spoofing and Adversarial Use}

The security use of fingerprints creates an incentive to manipulate them. Large commercial measurements have characterized both benign and adversarial browser fingerprints observed by production-facing services~\cite{wu2023manyfaces}. Gummy Browsers studies a targeted form of the threat: an attacker first collects a victim's browser fingerprint and then changes a separate browser through script injection, script modification, browser settings, or debugging tools so that fingerprint-linking systems associate the attacker with the victim~\cite{liu2022gummy}. The work demonstrates that browser-exposed surfaces are not immutable and motivates controlled manipulations of User-Agent, platform, WebGL, resources, and automation-facing properties.

HybridGuard adopts this adversarial premise but evaluates a narrower and more auditable question. The controlled-manipulation workstream records the exact tool, runner, version, configuration, and affected field paths, and provides verified baseline-to-active pairs. The main analysis does not infer an attack label from a tool name, a process return code, or the existence of a WebDriver marker. It first verifies that a catalogued field actually changed and only then asks whether a frozen relation becomes violated. This distinction is necessary because automation and debugging interfaces also have legitimate testing uses, and because a successfully executed intervention may not affect the fields observed by the current collector.

\subsection{Fingerprint Inconsistency and Configuration Relations}

The closest line of prior work detects contradictions inside browser fingerprints. FP-Scanner examines anti-fingerprinting countermeasures and shows that partial attribute modification can create combinations that do not naturally correspond to one browser or operating-system configuration~\cite{vastel2018fpscanner}. The scanner can use such contradictions to detect the presence of countermeasures and, for some attributes, infer the original environment. FP-Inconsistent revisits the idea in evasive bot traffic at larger scale. Its honey-site study shows that commercial bot services alter different fingerprint attributes to bypass anti-bot services, yet leave spatial inconsistencies within one fingerprint and temporal inconsistencies across observations~\cite{venugopalan2025fpinconsistent}. The paper learns rules from these inconsistencies and demonstrates that they expose part of the measured evasion.

BrowserFM attacks a related representation problem. It generates browser fingerprints under many Chromium versions, hardware environments, and configuration parameters, then links fingerprints, configurations, and environments through feature models and per-observation identifiers~\cite{huyghe2025browserfm}. The result supports queries such as which attributes respond to a particular switch, how headless and headful configurations differ, and which reduced attribute sets retain configuration information. Earlier cross-browser fingerprinting likewise demonstrates that OS- and hardware-level properties can remain observable across browser implementations~\cite{cao2017crossbrowser}, while GPU work shows that Web rendering can reveal low-level hardware variation~\cite{laor2022drawnapart}.

HybridGuard shares the premise that relationships carry information beyond isolated values. Its analysis boundary, however, is different. FP-Scanner and FP-Inconsistent primarily inspect browser-internal relations; BrowserFM studies configuration generation; cross-browser fingerprinting studies linkability. HybridGuard binds observations from Android Native APIs, the WebView host, and the in-app Web runtime to one collection provenance and evaluates explicit cross-layer semantics under controlled intervention. It also treats unavailable, unparseable, and not-applicable fields as distinct from matching values. Consequently, the closest browser systems are conceptual neighbors rather than drop-in baselines.

\subsection{Android WebView and Cross-Boundary Observability}

Android hybrid applications combine at least two programming and security domains. Native code controls the application process, permissions, package state, and WebView configuration; JavaScript executes under browser-origin and Web API semantics. Large-scale work on Android--Web hybridization documents the prevalence and architectural complexity of these combinations~\cite{tiwari2020hybridization}. Android's own documentation emphasizes that JavaScript bridges and WebView settings alter the boundary between untrusted Web content and application capabilities~\cite{android2026nativebridges,android2026webview,android2026websettings}.

Several systems analyze this boundary as an information-flow problem. BridgeTaint propagates taint information bi-directionally through JavaScript bridges and maps labels across the language boundary to identify privacy leaks and code injection paths~\cite{bai2019bridgetaint}. iWanDroid takes a demand-driven approach, selectively summarizing Native code according to how JavaScript consumes exposed functionality, so that information-flow violations can be followed in both directions~\cite{tiwari2023demand}. These works establish that a WebView cannot be analyzed faithfully by treating Native and JavaScript code as independent programs.

Recent measurement studies examine how the same boundary is used in deployed applications. Tracking Without Borders combines static and dynamic analysis over 13,045 Google Play applications and reports widespread WebView customization, browser fingerprinting, identifier bridging, cookie syncing, and access to permission-protected resources~\cite{weerasekara2025tracking}. Cross-Boundary Mobile Tracking introduces WebViewTracer, combining Frida-observed Java method calls with VisibleV8 JavaScript execution traces. On 10,000 applications, it identifies sensitive information injected from Java into WebViews and subsequently sent through Web APIs, while also observing fingerprinting behavior inside WebView scripts~\cite{datta2026crossboundary}. These results show that the host and Web runtime form an operationally coupled system even though their security contexts differ.

HybridGuard is adjacent to, but distinct from, this work. It does not attempt full source-to-sink taint tracking, SDK attribution, or detection of privacy exfiltration. Its inputs are fixed runtime observations collected by a purpose-built app, and its output is a versioned set of environment facts and relation states. A JavaScript bridge is therefore treated as host-topology evidence, not as proof of a leak; a WebView user-agent override may be legitimate and is encoded as a tolerance or not-assessed case rather than a universal alert.

\subsection{Detecting Fingerprinting Scripts versus Validating Environment Evidence}

Another body of work detects whether a Web page or application performs fingerprinting. FP-Fed trains privacy-preserving detectors for fingerprinting scripts~\cite{annamalai2024fpfed}; FP-tracer analyzes API-level information flow and entropy to identify fingerprinting behavior~\cite{boussaha2024fptracer}; Android-focused work detects fingerprinting activity inside application code~\cite{heid2024androidfp}. These systems answer a different question from HybridGuard. In HybridGuard, collection is intentional and declared. The question is whether the resulting Native, host, and Web observations form a provenance-complete and semantically compatible environment under the current experimental facts.

This task distinction changes both labels and evaluation. A fingerprinting-script detector may use script behavior as a positive label. HybridGuard must not: a page can legitimately read a fingerprint field, an automation interface can be used for testing, and a field may differ across containers without implying maliciousness. The project therefore separates collection metadata, verified intervention facts, model-visible evidence, and final relation status.

\subsection{Positioning Summary}

Table~\ref{tab:positioning} summarizes the closest task boundaries. The novelty claim in this draft is deliberately compositional: the contribution is the combination of Android multi-layer collection, per-field states, provenance-complete data units, explicit relation applicability, source-separated rules, and controlled two-state evaluation. The draft does not claim to be the first use of fingerprint consistency, the first cross-language WebView analysis, or a complete RBA system.

\begin{table*}[t]
\caption{Positioning against the closest lines of work. ``Controlled pairs'' means that the primary evaluation unit explicitly binds a baseline and a verified active intervention.}
\label{tab:positioning}
\small
\begin{tabularx}{\textwidth}{@{}p{2.8cm}p{2.5cm}p{2.4cm}p{2.2cm}X@{}}
\toprule
Line of work & Main observation boundary & Primary task & Controlled pairs & Difference from HybridGuard \\
\midrule
FP-Scanner / FP-Inconsistent & Browser attributes & Detect altered or evasive fingerprints & Not HybridGuard-style App pairs & Browser-internal spatial/temporal consistency; no Android Native--WebView provenance contract \\
Gummy Browsers & Victim and attacker browsers & Targeted fingerprint impersonation & Attack experiments, but different unit & Demonstrates spoofability; HybridGuard evaluates environment contradictions rather than identity linkage \\
BrowserFM & Browser, configuration, system & Structured configuration analysis & Configuration enumeration & Preserves generating configuration; does not evaluate Android cross-layer manipulation \\
BridgeTaint / iWanDroid & Java--JavaScript flows & Leakage, injection, information-flow violations & No & Code/data-flow analysis rather than runtime environment consistency \\
Tracking Without Borders / WebViewTracer & App, SDK, WebView, JavaScript, network & Privacy tracking and cross-context diffusion & No & Large-scale behavior measurement rather than fixed-contract evidence comparison \\
HybridGuard (current) & Android Native, WebView host, in-app Web & Provenance-aware relation violation & Yes: baseline $\rightarrow$ active & Reports field effects, applicability, alerts, and uncovered controlled configurations \\
HybridGuard (future) & Current layers plus external browser & Incremental Browser67 value & Planned & Requires sufficiently complete paired244 attack and benign data \\
\bottomrule
\end{tabularx}
\end{table*}
